\documentclass[a4paper]{article}     
\usepackage{amsfonts}
\usepackage{amsmath}
\usepackage{amssymb}
\usepackage{graphicx}
\usepackage{cancel}

\newcommand{\asa}{\Leftrightarrow} %als slechts als
\newcommand{\adR}{\Rightarrow} %als dan Rechts
\newcommand{\adL}{\Leftarrow}
\newcommand{\Lh}{\left(} %Links haakje
\newcommand{\Rh}{\right)}
\newcommand{\Lv}{\left[} %Links vierkanthaakje
\newcommand{\Rv}{\right]}
\newcommand{\La}{\left\{} %Links acolade
\newcommand{\Ra}{\right\}}
\newcommand{\Ras}{\right.} %Rechts acolade stelsel (omdat om stelsels te maken heb je een punt nodig
\newcommand{\pnR}{\rightarrow} %pijl naar Rechts
\newcommand{\limiet}{\lim\limits_}
\newcommand{\schrap}{\cancel}
\newcommand{\schrapb}{\bcancel} %backslash
\newcommand{\schrapk}{\xcancel} %kruis
\newcommand{\vervang}{\cancelto}
\newcommand{\intgrl}{\displaystyle\int} %integraal teken (bij het berekenen van primitieven)

\begin{document}

\noindent \large Auteur: Yoren Collier \\
\noindent \large Studierichting: Informatica\\
\noindent \large Oefeningenreeks 6B nr. 5.7\\

\medskip

\normalsize

$n$ als $x_3 = -36, x_n = -324$ en $q = -3$\\

$x_n = x_1 \cdot q^{n-1}$\\

$x_n = x_3 \cdot q^{n-3} \asa -324 = -36 \cdot \Lh -3 \Rh^{-3} \cdot q^n \asa -324 = \dfrac{-36}{\Lh -3 \Rh^3} \cdot q^n$\\

$\asa -324 = \dfrac{\schrap{-}36}{\schrap{-}27} \cdot q^n \asa -324 = \dfrac{4}{3} \cdot q^n \asa \Lh -3 \Rh^n = -324 \cdot \dfrac{3}{4}$\\

In geval dat je niet vanbuiten weet dat: $3^5 = 243$\\

$\asa \Lh -3 \Rh^n = -243 \asa \Lh -3 \Rh^n = -3 \cdot 81 \asa \Lh -3 \Rh^n = -3 \cdot \Lh -3 \Rh \cdot \Lh -27 \Rh \asa \Lh -3 \Rh^n = \Lh -3 \Rh^5 \asa n = 5$\\

\begin{thebibliography}{999}
%\bibitem{a} Referentie
\end{thebibliography}
\end{document} 